\section{Task 1 : Impedance Matching}
\label{sec:task1}

\subsection{Goals of this task}
The main objective of this step is to design a matching circuit to maximize energy transfert from
 an input impedance of an antenna and an output impedance of a Schottky diode. 
 In an RF harvesting system, an antenna circuit have an impedance ($Z_{in}$) of 50 $\Omega$. 
 This is a compromise between loss and power transmission. On the other hand, a Schottky diode 
 (who act like a charge) have a complex impédance ($Z_{diode}$) which is not fix regarding the frequency 
 of the system. Without any adaptation, there will be a high reflexion of catched energy by the antenna, 
 this problem will minimize the efficiency of our system.

\subsection{Modelling}

To simulate the behavior with the Keysight ADS software, we wanted to transform the complexe impédence 
of the diode to passives components like resistors, capacitors and inductors. There is the complex form of an impedance :  

\begin{equation} \label{eq:complex_form_impedance}
    Z = R + jX
\end{equation}

The real part ($R$) represent the resistive part of the diode. 
This is why this part is modelized by a simple resistor in the simulator.

The imaginary part ($X$) is the reactance of the component. 
With a Schottky diode, this reactance is negative, that means that we can replace this with a capacitor in the software

To modelize this imaginary part in the simulator, we need to calculate the capacitor’s value. 
We used this equation de determine this value

\begin{equation} \label{eq:capacitor_value}
    C = -\frac{1}{2\pi f X}
\end{equation}

\subsection{Expected results}

When the diode was modelized by its equivalent RC circuit, we added a LC filter between the source (50 $\Omega$) 
and the charge (Schottky diode). This filter’s objective is to compensate the imaginary part of the diode and 
keep a 5O $\Omega$ impedance at the expected frequency.

We use a method included in the Keysight’s software. This is a tool to tune dynamicaly the value of the 
inductance and the capacitor of the filter. The goal is to find a “dip” in the chart at the expected frequency. 
this “dip” is the optimal resonance of the system. Regarding the specifications, we want to obtain a $S_{11}$ 
parameter (reflexion coeficient) lower than -10 db at the wanted frequency.